\subsection{Operations on B-spaces}

Let B be a topological space.

Let X be a B-space and p its projection. Every topological subspace Y of X is endowed with the B-space structure whose projection is \(p \mid Y\) . Let A be a subspace of B; endowed with the map \(p_{\mathrm{A}}: \overline{p}^{1}(\mathrm{A}) \to \mathrm{A}\) induced by p by passing to subspaces, the topological space \(\overline{p}^{1}(\mathrm{A})\) is an A-space. It is called the A-space induced by \((\mathrm{X}, p)\) over A and is sometimes denoted \(X_{A}\) .

Let \((\mathrm{X}_{i})_{i\in\mathrm{I}}\) be a family of B-spaces, and let \(p_{i}\) be the projection of \(X_{i}\) . The sum space \(X = \coprod_{i\in I} X_{i}\) (TG, I, p. 15), equipped with the map \(p: X \to B\) defined by \(p(i,x) = p_{i}(x)\) (for \(i \in I\) and \(x \in X_{i}\) ), is a B-space called the sum of the family of B-spaces \((\mathrm{X}_{i})_{i\in\mathrm{I}}\) . The canonical injections \(X_{i} \to X\) are B-morphisms.

Let X be a B-space, p its projection, and let R be an equivalence relation on X. Denote by X/R the quotient space (TG, I, p. 20, def. 3). If the map \(p: X \rightarrow B\) is compatible with the relation R (E, II, p. 44), the map \(p'\colon X/R \to B\) induced by p by passing to the quotient is continuous (TG, I, p. 21, prop. 6); the B-space obtained by endowing X/R with the projection \(p'\) is then called the quotient B-space of X by the relation R.

